\section{总结}
我们证明了：只要所有发散都在规范不变的方案中正规化，准胶子算符的紫外发散实际上与准夸克算符的紫外发散相似。我们证明了全部 36 个纯准胶子算符的紫外发散都定域于时空之中，并且可以被乘法重正化而不相互混合，
\begin{align}
{\cal O}_{g}^{\mu \nu \rho\sigma}(\xi)=e^{-C_g |\xi_z|}Z_{wg}^{-1}Z_{vg1}^{-s/2}Z_{vg2}^{-(2-s)/2}{\cal O}_{bg}^{\mu \nu \rho\sigma}(\xi),
\end{align}
其中 $s$ 是为洛伦兹指标 $\{ \mu,\nu,\rho,\sigma\}$ 选出的 $z$ 分量个数，$C_g$、$Z_{wg}$、$Z_{vg1}$ 与 $Z_{vg2}$ 为重正化常数。
与准夸克算符一样，由于不相互混合，准胶子算符的强子矩阵元可以成为文献~\cite{Ma:2014jla,Ma:2017pxb}所定义的良好“格点截面”的例子：它们可在 LQCD 中计算并因子化为普通 PDF；通过对第一性原理 LQCD 计算产生的这些良好“格点截面”数据进行 QCD 全局分析，即可从中抽取 PDF。
